\chapter{Wightman Fields and Reconstruction}
\label{ch:volume-iv-wightman-fields-reconstruction}
\ConventionsLorentzian


\section{Field-Theoretic Data}

Let \(M=\mathbb M^D\) be \(D\)-dimensional Minkowski space with
translation coordinates \(x^\mu\) and metric
\(\eta=\operatorname{diag}(-1,1,\ldots,1)\).  The proper orthochronous
Poincare group is denoted by
\[
  \Poinc=\mathbb R^D\rtimes SO^+(1,D-1).
\]
The Wightman framework describes a QFT by Hilbert-space data together with
local operator-valued distributions.  The basic data are:
\[
  (\Hilb,\vac,U,\mathcal D,\{\widehat\Phi_\alpha\}_{\alpha\in I}).
\]
Here \(\Hilb\) is a Hilbert space, \(\vac\in\Hilb\) is a unit vector, \(U\) is a
strongly continuous unitary representation of \(\Poinc\),
\(\mathcal D\subset\Hilb\) is a dense invariant domain, \(I\) is the field-label
index set, and each \(\widehat\Phi_\alpha\) is
an operator-valued tempered distribution on \(M\) with common domain
\(\mathcal D\).  Thus this chapter uses the Schwartz-space branch of the
test-function convention fixed in Chapter~\ref{ch:starting-data-local-qft};
compactly supported local smearings are obtained by restriction.  For a
Schwartz test function \(f\in\mathcal S(M)\),
\[
  \widehat\Phi_\alpha(f)=\int_M \dd^D x\, f(x)\widehat\Phi_\alpha(x)
\]
denotes the smeared operator.  This integral notation means distributional
pairing, and all products of fields below are products of smeared operators on
the common domain.

The field labels carry an antilinear involution.  Let
\[
  E_I=\bigoplus_{\alpha\in I}\mathbb C e_\alpha
\]
be the finite field-label space used in a given multiplet, and let
\[
  J:E_I\to E_I,\qquad J e_\alpha=e_{\alpha^\dagger},
  \qquad J^2=1,
\]
be an antilinear map; if the chosen field basis is not closed under componentwise
conjugation, \(J\) includes the corresponding finite-dimensional conjugation
matrix.  For an \(E_I\)-valued test function \(h=\sum_\alpha h^\alpha e_\alpha\)
define
\[
  \widehat\Phi(h)=\sum_\alpha \widehat\Phi_\alpha(h^\alpha),
  \qquad
  h^\dagger(x)=Jh(x)
  =
  \sum_\alpha \overline{h^\alpha(x)}\,e_{\alpha^\dagger}.
\]
The Hermiticity hypothesis for the operator-valued distribution is the domain
inclusion
\begin{equation}
  \widehat\Phi(h)^*\supset \widehat\Phi(h^\dagger).
  \label{eq:wightman-field-adjoint-inclusion}
\end{equation}
Equivalently, for a scalar test function \(f\),
\[
  \widehat\Phi_\alpha(f)^*
  \supset
  \widehat\Phi_{\alpha^\dagger}(\overline f).
\]
Here \(\widehat\Phi(h)^*\) is the Hilbert-space adjoint of the densely defined
operator \(\widehat\Phi(h):\mathcal D\to\mathcal D\); the inclusion means that
\(\mathcal D\subset\operatorname{Dom}(\widehat\Phi(h)^*)\) and that the two
operators agree on \(\mathcal D\).  Thus a Hermitian field smeared with a real
test function is symmetric on \(\mathcal D\).  Essential self-adjointness of its
closure is an additional analytic assertion, not part of the formal adjoint
identity.

The vacuum is invariant:
\[
  U(g)\vac=\vac,\qquad g\in \Poinc .
\]
For a pure translation we use the abbreviation \(U(a):=U(a,1)\).  Translations
are written
\[
  U(a)=U(a,1)=\exp(\ii a^\mu P_\mu),
  \qquad a\in\mathbb R^D,
\]
in the convention of Chapter \ref{ch:starting-data-local-qft}.  Let \(E_P\) be
the joint spectral measure of \(P=(P_0,\ldots,P_{D-1})\).  With
\(p^2=\eta_{\mu\nu}p^\mu p^\nu\), define the open and closed forward cones by
\[
  V_+
  =
  \{p\in\mathbb R^D\mid p^0>0,\ p^2<0\},
  \qquad
  \overline V_+
  =
  \{p\in\mathbb R^D\mid p^0\ge0,\ p^2\le0\}.
\]
The spectrum condition is
\[
  \operatorname{Spec}(P)\subset \overline V_+ .
\]
If the field multiplet transforms in a finite-dimensional Lorentz
representation \(R\), covariance is the identity
\[
  U(a,\Lambda)^{-1}\widehat\Phi_\alpha(f)U(a,\Lambda)
  =
  R(\Lambda)_\alpha{}^\beta\,
  \widehat\Phi_\beta(f_{a,\Lambda}),
\]
where
\[
  f_{a,\Lambda}(x)=f\bigl(\Lambda^{-1}(x-a)\bigr).
\]
Equivalently, for a pure translation, conjugating in the opposite direction
gives
\[
  U(a)\widehat\Phi_\alpha(f)U(a)^{-1}
  =
  \widehat\Phi_\alpha(f^{(+a)}),
  \qquad
  f^{(+a)}(x)=f(x+a).
\]
The precise spin and internal indices are carried by the label \(\alpha\).

Locality is imposed after smearing.  If the supports of \(f\) and \(g\) are
spacelike separated, then
\[
  [\widehat\Phi_\alpha(f),\widehat\Phi_\beta(g)]_{\rm gr}=0
\]
on \(\mathcal D\), where
\[
  [A,B]_{\rm gr}=AB-(-1)^{|A||B|}BA
\]
for homogeneous field parities \(|A|,|B|\in\mathbb Z/2\mathbb Z\).  For purely
bosonic fields this is the ordinary commutator; when both fields are odd it is
the anticommutator.

When Wightman fields are given as operator data, cyclicity is part of the
chosen field presentation: the Hilbert space is the vacuum representation
generated by the specified fields.  When the starting point is instead a
correlation hierarchy, the reconstruction theorem below constructs this cyclic
representation from the hierarchy itself.

The vacuum is cyclic for the polynomial field algebra:
\[
  \operatorname{span}\{
    \widehat\Phi_{\alpha_1}(f_1)\cdots\widehat\Phi_{\alpha_n}(f_n)\vac
    \mid n\ge0
  \}
  \quad\text{is dense in }\Hilb .
\]

\begin{figure}[H]
  \centering
  \resizebox{0.90\textwidth}{!}{%
  \begin{tikzpicture}[x=1cm,y=1cm,every node/.style={font=\scriptsize}]
    \tikzset{
      box/.style={draw,rounded corners=4pt,line width=0.8pt,fill=blue!4,
        minimum width=3.0cm,minimum height=0.90cm,align=center},
      arrow/.style={-{Latex[length=2mm]},line width=0.85pt},
      note/.style={align=center}
    }
    \node[box] (hilb) at (-4.2,1.3) {Hilbert space\\\(\Hilb,\vac,U\)};
    \node[box] (domain) at (0,1.3) {common domain\\\(\mathcal D\)};
    \node[box] (fields) at (4.2,1.3) {smeared fields\\\(\widehat\Phi_\alpha(f)\)};
    \node[box] (spectrum) at (-3.2,-1.0)
      {spectrum\\\(\operatorname{Spec}(P)\subset\overline V_+\)};
    \node[box] (locality) at (3.2,-1.0)
      {locality\\spacelike commutators};
    \draw[arrow] (hilb)--(domain);
    \draw[arrow] (domain)--(fields);
    \draw[arrow] (hilb)--(spectrum);
    \draw[arrow] (fields)--(locality);
    \draw[arrow] (locality)--(spectrum);
    \node[align=center] at (0,-1.72) {analytic structure};
    \node[note,green!45!black] at (0,-2.35)
      {The field data are distributional, domain-sensitive, and local after smearing.};
  \end{tikzpicture}
  }
  \caption{The Wightman framework packages local fields, a Hilbert space,
  covariance, spectral positivity, and locality into one distributional
  structure.}
  \label{fig:wightman-data}
\end{figure}

\section{Fields and Local Observable Content}

The objects \(\widehat\Phi_\alpha(f)\) in the Wightman data are generally unbounded
operators.  Their algebraic products are defined first on the common invariant
domain \(\mathcal D\), and their locality statement is a statement about
graded commutators on that domain.  A local algebra of bounded observables is
a further object constructed from, or compared with, these field
coordinatizations.

For an open region \(\mathcal O\subset M\), let
\[
  \mathfrak P_\Phi(\mathcal O)
\]
denote the unital \(*\)-algebra on \(\mathcal D\) generated by smeared fields
\(\widehat\Phi_\alpha(f)\) with \(\operatorname{supp}f\subset\mathcal O\), together
with their adjoints on \(\mathcal D\).  The subscript \(\Phi\) names the chosen
field family; the individual Hilbert-space fields in the family remain the
hatted operator-valued distributions \(\widehat\Phi_\alpha\).  This is an
algebra of generally unbounded operators on a common domain, not a
\(C^*\)-algebra.  If the
relevant smeared Hermitian fields are essentially self-adjoint and if the
closures and bounded Borel functions have the required localization
properties, one can form a represented local von Neumann algebra
\begin{equation}
  \mathcal R_\Phi(\mathcal O)
  =
  \{\text{bounded functions of closures of fields localized in }\mathcal O\}'',
  \label{eq:wightman-generated-local-vn-algebra}
\end{equation}
or, more invariantly, require the unbounded fields to be affiliated with a
local algebra \(\mathcal R(\mathcal O)\).  The double prime denotes weak
operator closure.  The affiliation formulation is often the correct one:
spectral projections of a self-adjoint operator affiliated with
\(\mathcal R(\mathcal O)\) belong to \(\mathcal R(\mathcal O)\).

\begin{figure}[H]
  \centering
  \resizebox{0.96\textwidth}{!}{%
  \begin{tikzpicture}[x=1cm,y=1cm,every node/.style={font=\scriptsize}]
    \tikzset{
      box/.style={draw,rounded corners=4pt,line width=0.8pt,fill=blue!4,
        text width=2.85cm,minimum height=0.95cm,inner sep=4pt,align=center},
      obox/.style={draw,rounded corners=4pt,line width=0.8pt,fill=orange!8,
        text width=2.85cm,minimum height=0.95cm,inner sep=4pt,align=center},
      arrow/.style={-{Latex[length=2mm]},line width=0.85pt},
      darrow/.style={-{Latex[length=2mm]},dashed,line width=0.85pt}
    }
    \node[box] (fields) at (-5.4,0.65)
      {field coordinates\\\(\widehat\Phi_\alpha(f)\)};
    \node[box] (poly) at (-1.8,0.65)
      {domain algebra\\\(\mathfrak P_\Phi(\mathcal O)\)};
    \node[obox] (aff) at (1.8,0.65)
      {affiliation\\or bounded\\functional calculus};
    \node[obox] (net) at (5.4,0.65)
      {local observable\\algebra\\\(\mathcal R(\mathcal O)\)};
    \node[box,fill=green!7] (other) at (-1.8,-1.45)
      {another field\\coordinatization};
    \draw[arrow] (fields.east)--(poly.west);
    \draw[arrow] (poly.east)--(aff.west);
    \draw[arrow] (aff.east)--(net.west);
    \draw[darrow] (other.east) to[out=0,in=-120] (net.south);
    \node[align=center,fill=white,inner sep=1.5pt] at (3.15,-1.78)
      {same local content\\if proven};
    \node[align=center] at (1.8,1.55)
      {extra domain and closure hypotheses};
  \end{tikzpicture}
  }
  \caption{Wightman fields are distributional coordinates.  A local observable
  algebra is obtained only after specifying how the unbounded smeared fields
  are affiliated with bounded local algebras, or after constructing the net
  directly.}
  \label{fig:wightman-fields-to-local-algebras}
\end{figure}

Several different field families may coordinatize the same local observable
content.  In this monograph a field coordinatization is treated as part of the
presentation, while the local observable net, when it has been constructed, is
the invariant local object.  Thus the Wightman hierarchy for a chosen field
family is complete for reconstructing that chosen field presentation, whereas
comparisons between different field presentations require an additional
equivalence statement: one must show that they generate the same represented
local algebras, or a specified subtheory, in the relevant vacuum sector.

\section{Wightman Functions}

The vacuum \(n\)-point distributions are
\begin{equation}
  W_{\alpha_1\cdots \alpha_n}(x_1,\ldots,x_n)
  =
  \bra{\vac}
  \widehat\Phi_{\alpha_1}(x_1)\cdots\widehat\Phi_{\alpha_n}(x_n)
  \ket{\vac}.
  \label{eq:wightman-functions}
\end{equation}
They are tempered distributions on \(M^n\).  Translation invariance allows one
to pass to difference variables
\[
  \xi_j=x_{j+1}-x_j,\qquad j=1,\ldots,n-1.
\]
The spectral condition implies that the Fourier transform in the difference
variables has support in \(\overline V_+^{\,n-1}\).  This is the same positive
spectral input that will lead in Chapter \ref{ch:kallen-lehmann} to the
K\"all\'en--Lehmann representation of two-point functions.

The Wightman hierarchy also records the adjunction operation on fields.  On
one-field test functions the adjunction is the map \(h\mapsto h^\dagger\) fixed
in \eqref{eq:wightman-field-adjoint-inclusion}.  On test functions valued in
field-label tensor powers, one must also reverse the order of insertions because
the Hilbert-space adjoint reverses products.  For
\(f_n\in\mathcal S(M^n,E_I^{\otimes n})\), define the antilinear tensor
reversal \(R_n:E_I^{\otimes n}\to E_I^{\otimes n}\) by
\[
  R_n(v_1\otimes\cdots\otimes v_n)
  =
  Jv_n\otimes\cdots\otimes Jv_1
\]
and set
\[
  f_n^\dagger(x_1,\ldots,x_n)
  =
  R_n\bigl(f_n(x_n,\ldots,x_1)\bigr).
\]
Thus \(f_n^\dagger\) simultaneously reverses the order of insertions, complex
conjugates the coefficient functions, and applies the field-label adjunction to
each tensor factor.

The Wightman hierarchy satisfies five structural conditions that are directly
read from the Hilbert-space data:
\begin{description}[leftmargin=2.9cm,style=nextline]
\item[Covariance.]
The distribution \(W_n\) transforms under \(\Poinc\) by the
\(n\)-fold field representation on its labels and test functions.
\item[Spectrum.]
In translation-difference variables,
\[
  \operatorname{supp}\widehat W_n\subset \overline V_+^{\,n-1}.
\]
\item[Hermiticity.]
For all labels and test functions,
\[
  \overline{W_{\alpha_1\cdots\alpha_n}
    (f_1\otimes\cdots\otimes f_n)}
  =
  W_{\alpha_n^\dagger\cdots\alpha_1^\dagger}
  (\overline f_n\otimes\cdots\otimes\overline f_1).
\]
This is the vacuum distribution form of the operator-domain adjoint condition
\(\widehat\Phi(h)^*\supset\widehat\Phi(h^\dagger)\), with product order
reversed after taking adjoints.
\item[Locality.]
\(W_n\) has the graded permutation property when adjacent field insertions
are smeared with spacelike separated test functions.
\item[Positivity.]
For every finite sequence \(F=(f_0,f_1,f_2,\ldots)\) of test functions,
\[
  \sum_{m,n} W_{m+n}(f_m^\dagger\otimes f_n)\ge0 .
\]
\end{description}
In the positivity condition, \(f_n\) denotes a compactly supported or Schwartz
test function valued in the appropriate finite direct sum of field labels.
Hermiticity is the compatibility condition that makes the
reconstructed left-multiplication fields have the intended adjoints on the
common domain.

\section{Clustering and the Vacuum Sector}

The vacuum vector should represent one pure infinite-volume phase.  In the
Wightman setting this is expressed by a cluster property.  For two polynomial
local operators \(A\) and \(B\), formed from fields smeared in bounded
regions, define
\[
  B(a)=U(a)BU(a)^{-1}.
\]
The weak cluster property in the vacuum sector is
\begin{equation}
  \lim_{\lambda\to+\infty}
  \langle\vac,A\,B(\lambda a)\vac\rangle
  =
  \langle\vac,A\vac\rangle
  \langle\vac,B\vac\rangle
  \label{eq:wightman-weak-cluster}
\end{equation}
for every spacelike vector \(a\).  In terms of Wightman functions, this says
that widely spacelike-separated groups of insertions factorize.

The spectral argument is as follows.  Insert the spectral resolution of
translations between \(A\) and \(B(a)\).  The vacuum atom at \(p=0\) gives the
product on the right side of \eqref{eq:wightman-weak-cluster}.  If the vacuum
is the only translation-invariant vector, then no other atom at \(p=0\)
contributes.  The remaining continuous spectral part has oscillatory phase
\(\exp(\ii\lambda a\cdot p)\); after smearing, its contribution vanishes in
the weak large-separation limit.  With a mass gap and suitable regularity, the
same argument gives stronger decay estimates.  Selecting a pure vacuum sector
is the condition that the zero-momentum projection is the one-dimensional
projection onto \(\mathbb C\vac\).

\begin{theorem}[Cluster decomposition and vacuum uniqueness]
\label{thm:wightman-cluster-vacuum-uniqueness}
Let
\((\Hilb,\vac,U,\mathcal D,\{\widehat\Phi_\alpha\})\)
be Wightman field data satisfying the assumptions above, and let
\[
  P_0=E(\{0\})
\]
be the joint spectral projection of the translation generators onto zero
momentum.  The following conditions are equivalent.
\begin{enumerate}[label=\textup{(\roman*)}]
\item The zero-momentum subspace is one-dimensional:
\[
  P_0\Hilb=\mathbb C\vac .
\]
\item The vacuum is the unique translation-invariant vector up to scalar
multiple.
\item The weak cluster property \eqref{eq:wightman-weak-cluster} holds for
all polynomial field insertions \(A\) and \(B\) and for every spacelike
translation direction \(a\).
\end{enumerate}
\end{theorem}

\begin{proof}
The equivalence of (i) and (ii) is the joint spectral theorem for the
translation representation: translation-invariant vectors are precisely the
vectors in \(\operatorname{Ran}P_0\).  For (i) \(\Rightarrow\) (iii), write
\[
  \langle\vac,A\,B(\lambda a)\vac\rangle
  =
  \langle A^\dagger\vac,U(\lambda a)B\vac\rangle ,
\]
first for polynomial fields whose domain adjoints are defined on
\(\mathcal D\), and then by linearity.  The Wightman spectral condition and
temperedness make this matrix element the Fourier transform, in the weak
large-spacelike-separation sense, of the spectral distribution between
\(A^\dagger\vac\) and \(B\vac\).  Its zero-momentum atom contributes
\[
  \langle A^\dagger\vac,P_0 B\vac\rangle
  =
  \langle\vac,A\vac\rangle\,\langle\vac,B\vac\rangle ,
\]
because \(P_0=|\vac\rangle\langle\vac|\).  The remaining spectral part has no
zero-momentum atom.  The analytic input is the following consequence of the
Wightman assumptions.  For vectors of the form
\(\psi=A^\dagger\vac\) and \(\chi=B\vac\) with \(P_0\chi=0\), the matrix
coefficient \(\langle\psi,U(\lambda a)\chi\rangle\) tends to zero as
\(\lambda\to+\infty\) for spacelike \(a\).  Indeed, after smearing the field
insertions, this coefficient is the boundary value of a tempered analytic
function in the forward tube.  Locality identifies the boundary values obtained
from the two spacelike orderings at the corresponding Jost points.  The
edge-of-the-wedge continuation removes singular support at finite nonzero
spacelike separation; the only translation-invariant spectral atom compatible
with all spacelike translates is the atom at \(p=0\).  After subtracting
\(P_0\), the remaining one-parameter spectral distribution has locally
integrable density in the separation variable relevant to the smeared matrix
element, and its Fourier transform vanishes at infinity by the
Riemann--Lebesgue lemma.  This gives the required weak large-separation limit.

Conversely, assume the weak cluster property.  The same spectral decomposition
shows that the large-separation limit of the non-oscillatory part is
\(\langle A^\dagger\vac,P_0B\vac\rangle\).  Comparing with
\eqref{eq:wightman-weak-cluster} gives
\[
  \langle A^\dagger\vac,P_0B\vac\rangle
  =
  \langle A^\dagger\vac,\vac\rangle\,\langle\vac,B\vac\rangle
\]
for all polynomial \(A,B\).  Since the polynomial field orbit of \(\vac\) is
dense, this identity implies
\[
  P_0B\vac=\vac\,\langle\vac,B\vac\rangle
\]
on a dense set of vectors \(B\vac\).  Hence \(P_0=|\vac\rangle\langle\vac|\),
which is (i).
\end{proof}

\section{Reconstruction From Correlation Functions}

The Wightman reconstruction theorem says that the correlation hierarchy is
itself a complete encoding of a chosen Wightman field presentation, provided
the hierarchy satisfies the Wightman conditions with the required regularity.

\begin{theorem}[Wightman reconstruction]
Let \(\{W_n\}_{n\ge0}\) be a hierarchy with \(W_0=1\), where \(W_n\) is a
continuous linear functional on
\(\mathcal S(M^n,E_I^{\otimes n})\).  Equivalently, on every finite
field-label subspace, each component distribution is bounded by finitely many
Schwartz seminorms: after choosing scalar components \(F_j\) of a test function
\(F\), there are \(C,N<\infty\) such that
\[
  |W_n(F)|
  \le
  C\sum_j\sum_{|a|,|b|\le N}
  \sup_{x\in M^n}|x^a\partial^b F_j(x)| .
\]
Suppose the hierarchy satisfies Poincare covariance, the spectral support
condition, Hermiticity, locality, and positivity.  Then there exists Wightman
field data
\[
  (\Hilb,\vac,U,\mathcal D,\{\widehat\Phi_\alpha\}_{\alpha\in I})
\]
with the following properties:
\begin{enumerate}[label=\textup{(\roman*)}]
\item Its vacuum correlation functions are exactly \(\{W_n\}\).
\item The dense domain is the polynomial field orbit of the vacuum,
\[
  \mathcal D
  =
  \operatorname{span}\{
    \widehat\Phi_{\alpha_1}(f_1)\cdots
    \widehat\Phi_{\alpha_n}(f_n)\vac
    \mid n\ge0,\ \alpha_j\in I,\ f_j\in\mathcal S(M)
  \}.
\]
In particular, \(\vac\) is cyclic for the reconstructed polynomial field
algebra.
\item \(U\) is a strongly continuous unitary representation of \(\Poinc\), the
fields are covariant operator-valued tempered distributions on the common
domain \(\mathcal D\), and the spectrum, adjoint, and locality properties
specified above hold on \(\mathcal D\).
\end{enumerate}
The cyclic reconstruction satisfying these properties is unique up to the
canonical unitary induced by identifying finite test-function sequences with
the same Wightman inner products.

\end{theorem}

For a hierarchy satisfying the reconstruction hypotheses, the reconstructed
translation representation has a zero-momentum projection \(P_0\).  Adding the
weak cluster condition \eqref{eq:wightman-weak-cluster} selects exactly the
case \(P_0\Hilb=\mathbb C\vac\) by Theorem
\ref{thm:wightman-cluster-vacuum-uniqueness}.  Thus clustering is the
correlation-function condition for reconstructing a pure vacuum sector.

The uniqueness statement is internal to the specified hierarchy and its field
labels.  It says that two cyclic reconstructions of the same \(\{W_n\}\) have
the same pre-Hilbert space of finite test-function sequences modulo null
vectors, hence the same vacuum matrix elements for every finite product of
field insertions.  This is a uniqueness statement about the Wightman hierarchy
and the cyclic field presentation it generates.  It does not identify different
choices of field coordinatization whose local observable algebras may agree
only after additional closure, affiliation, or quotient constructions.  Those
comparisons belong to the relation between Wightman fields and local nets
described above and in Chapter
\ref{ch:volume-iv-aqft-local-nets-states}.

The construction begins with formal vectors
\[
  [\alpha_1,f_1;\ldots ;\alpha_n,f_n]
\]
and defines an inner product by the Wightman distributions.  Positivity makes
this a positive semidefinite form.  Dividing by zero-norm vectors and
completing gives \(\Hilb\).  Fields act by left multiplication on the dense
subspace of finite sequences, and the Poincare action is induced by covariance
of the hierarchy.

More explicitly, let \(\mathcal V\) be the vector space of finite sequences
\[
  F=(f_0,f_1,f_2,\ldots),
\]
where \(f_n\) is a Schwartz test function on \(M^n\) valued in the \(n\)-fold
field-label space and only finitely many components are nonzero.  Define
\begin{equation}
  \langle F,G\rangle_W
  =
  \sum_{m,n}
  W_{m+n}(f_m^\dagger\otimes g_n).
  \label{eq:wightman-reconstruction-inner-product}
\end{equation}
Let
\[
  \mathcal N_W=\{F\in\mathcal V\mid \langle F,F\rangle_W=0\}.
\]
The dense domain \(\mathcal D\) is the image of \(\mathcal V/\mathcal N_W\)
inside its Hilbert-space completion.  For an \(E_I\)-valued one-field test
function \(h\), the reconstructed smeared field is
\[
  \widehat\Phi(h)[F]=[h\otimes F],
\]
where \(h\otimes F\) means left insertion of the one-field test function before
the sequence \(F\).  In components,
\(\widehat\Phi_\alpha(k)[F]=[(k e_\alpha)\otimes F]\).  Positivity and the
Wightman identities
ensure that this action is well defined on the null quotient and that the
vacuum vector is the class of \((1,0,0,\ldots)\).  The Poincare
representation acts by transforming all test functions; covariance of
\(\{W_n\}\) makes this action unitary.

The same finite-sequence construction gives cyclicity.  Starting with the
vacuum class \(\vac=[(1,0,0,\ldots)]\), repeated left insertion gives
\[
  \widehat\Phi(h_1)\cdots\widehat\Phi(h_n)\vac
  =
  [h_1\otimes\cdots\otimes h_n],
\]
and finite sums of such vectors are precisely the image of \(\mathcal V\) in
the quotient.  Since \(\mathcal D\) is defined as this image, \(\mathcal D\) is
the polynomial field orbit of \(\vac\), and its Hilbert-space closure is all of
\(\Hilb\).

The same formula verifies the adjoint inclusion on the reconstructed common
domain.  For an \(E_I\)-valued one-field test function \(h\) and finite
sequences \(F,G\),
\[
  \langle \widehat\Phi(h)F,G\rangle_W
  =
  \langle F,\widehat\Phi(h^\dagger)G\rangle_W .
\]
Indeed, inserting \(h\) on the left of \(F\) and then applying the dagger in
\eqref{eq:wightman-reconstruction-inner-product} moves \(h\) to the right
place with the label adjunction and complex conjugation fixed above.  After
quotienting by \(\mathcal N_W\), this identity says precisely that
\[
  \widehat\Phi(h)^*\supset \widehat\Phi(h^\dagger)
\]
on the dense finite-sequence domain.

\begin{figure}[H]
  \centering
  \resizebox{0.96\textwidth}{!}{%
  \begin{tikzpicture}[x=1cm,y=1cm,every node/.style={font=\scriptsize}]
    \tikzset{
      box/.style={draw,rounded corners=4pt,line width=0.8pt,fill=orange!7,
        text width=2.85cm,minimum height=0.95cm,inner sep=4pt,align=center},
      arrow/.style={-{Latex[length=2mm]},line width=0.85pt}
    }
    \node[box] (hier) at (-5.8,0) {distribution hierarchy\\\(\{W_n\}\)};
    \node[box] (pre) at (-1.75,0) {finite sequences\\of test functions};
    \node[box] (quot) at (1.75,0) {quotient by\\null vectors};
    \node[box] (hilb) at (5.55,0) {completion\\\(\Hilb\)};
    \node[box,fill=green!7] (ops) at (0,-1.85) {fields by\\left multiplication};
    \draw[arrow] (hier.east)--(pre.west)
      node[midway,above=12pt,fill=white,inner sep=1.2pt] {positivity};
    \draw[arrow] (pre.east)--(quot.west);
    \draw[arrow] (quot.east)--(hilb.west);
    \draw[arrow] (pre.south)--(ops.north west);
    \draw[arrow] (ops.north east)--(hilb.south west);
  \end{tikzpicture}
  }
  \caption{Wightman reconstruction builds the Hilbert space and fields from the
  positive correlation hierarchy.}
  \label{fig:wightman-reconstruction}
\end{figure}

\section{Analytic Continuation From Spectral Positivity}

For scalar fields, write the Wightman function in difference variables as
\[
  W_n(\xi_1,\ldots,\xi_{n-1}).
\]
The support condition for the Fourier transform implies that \(W_n\) is the
boundary value of a holomorphic function on the forward tube
\[
  T_{n-1}
  =
  \{z_j=\xi_j+\ii\eta_j\in\mathbb C^D
  \mid \eta_j\in V_+,\ j=1,\ldots,n-1\}.
\]
Lorentz covariance enlarges the domain to the extended tube obtained by the
complex Lorentz group.  Locality then identifies the analytic continuations
associated with permutations at Jost configurations, where all relevant real
linear combinations are spacelike.

\begin{figure}[H]
  \centering
  \resizebox{0.78\textwidth}{!}{%
  \begin{tikzpicture}[x=1cm,y=1cm,every node/.style={font=\scriptsize}]
    \tikzset{
      cone/.style={draw,fill=blue!5,line width=0.85pt},
      arrow/.style={-{Latex[length=2mm]},line width=0.85pt},
      box/.style={draw,rounded corners=4pt,fill=green!7,line width=0.8pt,
        minimum width=2.75cm,minimum height=0.80cm,align=center}
    }
    \draw[->] (-4.4,0)--(4.4,0) node[right] {\(\operatorname{Re}z\)};
    \draw[->] (0,-0.35)--(0,3.0) node[above] {\(\operatorname{Im}z\)};
    \draw[cone] (0,0)--(-1.3,2.45)--(1.3,2.45)--cycle;
    \node[blue!55!black] at (0,1.45) {\(V_+\)};
    \node[box] (tube) at (-3.0,2.25) {forward tube\\\(M+\ii V_+\)};
    \node[box] (ext) at (3.0,2.25) {extended tube\\complex Lorentz orbit};
    \node[box] (jost) at (0,-0.95) {Jost points\\real spacelike configurations};
    \draw[arrow] (tube)--(ext) node[midway,above] {covariance};
    \draw[arrow] (jost)--(tube) node[midway,left] {boundary values};
    \draw[arrow] (jost)--(ext) node[midway,right] {locality};
  \end{tikzpicture}
  }
  \caption{Spectral support gives tube analyticity; covariance and locality
  enlarge and glue the analytic domains.}
  \label{fig:wightman-tube-analyticity}
\end{figure}

The Euclidean Schwinger functions are restrictions of these analytic functions
to noncoincident imaginary-time configurations.  The converse passage from
Euclidean data to Lorentzian Wightman data requires reflection positivity and
growth hypotheses.  That converse is the subject of the next chapter.

\section{Observable Subpresentations in Gauge Theories}

A gauge construction enters the positive Hilbert-space Wightman framework
through a specified observable sector.  The object to verify is the following
positive local field presentation.

\begin{definition}[observable Wightman subpresentation]
Suppose a gauge construction supplies a positive physical Hilbert space
\(\Hilb_{\rm phys}\), a vacuum vector \(\vac_{\rm phys}\), a unitary Poincare
representation \(U_{\rm phys}\), a common dense domain
\(\mathcal D_{\rm phys}\), and a local net of physical von Neumann algebras
\(\mathcal R_{\rm phys}(\mathcal O)\).  A family of gauge-invariant local
distributions \(\{O_\gamma\}_{\gamma\in J}\) is an observable Wightman
subpresentation if:
\begin{enumerate}[label=(\roman*)]
\item each \(O_\gamma(f)\), \(f\in\mathcal S(M)\), is an operator on
\(\mathcal D_{\rm phys}\) and the map \(f\mapsto O_\gamma(f)\) is a tempered
operator-valued distribution;
\item the vacuum, covariance, spectrum, adjunction, locality, and positivity
conditions of the Wightman framework hold on \(\Hilb_{\rm phys}\);
\item for every open region \(\mathcal O\), all \(O_\gamma(f)\) with
\(\operatorname{supp}f\subset\mathcal O\) are affiliated with
\(\mathcal R_{\rm phys}(\mathcal O)\), and the local von Neumann algebra
generated by their bounded functional calculi is the specified observable
subnet under discussion.
\end{enumerate}
\end{definition}

In a BRST realization, the auxiliary space \(\mathcal K\) carries an odd
nilpotent operator \(Q\) and the physical Hilbert space is obtained, after the
required completion and quotient by null vectors, from
\[
  \Hilb_{\rm phys}
  =
  \ker Q\,/\,\overline{\operatorname{im}Q}.
\]
A local auxiliary operator \(B\) descends to a physical local observable when it
preserves \(\ker Q\) and \(\overline{\operatorname{im}Q}\); a sufficient
algebraic condition on a common invariant domain is
\[
  [Q,B]_{\rm gr}=0.
\]
Operators differing by a graded commutator \([Q,C]_{\rm gr}\) induce the same
operator on BRST cohomology, provided the same domain and closure hypotheses
hold.  The Wightman functions to which the reconstruction theorem applies are
therefore the vacuum matrix elements of the descended observable fields
\(\{O_\gamma\}\) on \(\Hilb_{\rm phys}\); auxiliary representatives enter
through their cohomology classes.

Thus the Wightman framework supplies precise theorems for local observable
fields, spectral positivity, analytic continuation, and reconstruction.  Its
output is Hilbert-space QFT data with operator-valued distributions and a
distinguished vacuum representation; comparison between gauges is a comparison
of the resulting observable nets or observable Wightman subpresentations.
